\section{The Geometry of Central-Path Conditioning}
\label{sec:conditioning}

\subsection{Question, answer, and relation to prior work}
\label{sec:conditioning-overview}

Quantum interior-point methods that invoke a quantum linear-system solver, and
classical interior-point methods that solve Newton equations inexactly, pay for
the conditioning of a Newton matrix.  A familiar estimate is
$\kappa=\Theta(\mu^{-2})$, but the classical derivations concern the
logarithmic barrier and often a primal--dual Schur complement rather than the
reduced primal Hessian.  It is therefore natural to ask whether a different
self-concordant barrier can remove the ill-conditioning.

For reduced barrier Hessians the answer is no in the relevant minimax sense.
The obstruction is the Euclidean geometry of the objective sublevel set.  We
prove a lower bound for every self-concordant barrier, pin every eigenvalue
between directional chord widths, and show that the canonical log and log-det
barriers attain the condition-number floor up to an instance-dependent
constant.  Error-bound exponents then give a dictionary of conditioning laws.

The universal quantifier over barriers is important.  Earlier conditioning
analyses concern the log barrier, including Wright's analysis of logarithmic
barrier Hessians and the survey of Forsgren, Gill, and Wright
\cite{wright1994-properties-hessian-logarithmic-barrier,forsgren2002-interior-methods-nonlinear-optimization}.
The $\mu^{-2}$-type estimates in
\citet[Section~7.4]{augustino2023-quantum-interior-point-methods-semidefinite},
for example, are for primal--dual Schur systems.  Conversely,
\citet{allamigeon2022-no-self-concordant-barrier} prove a lower bound uniform
over self-concordant barriers, but for the number and geometry of central-path
iterations, not Hessian conditioning.  The error-bound consequences below
connect to weak sharp minima \cite{burke1993-weak-sharp-minima} and to Sturm's
singularity-degree bounds for semidefinite systems
\cite{sturm2000-error-bounds-linear-matrix-inequalities}.

\begin{remark}[Novelty and textbook ingredients]
\label{rem:conditioning-novelty}
The Dikin-ellipsoid and asymmetric-containment facts used in the proofs, as
well as the centrality identity
$\Pi_{\mathcal V}\nabla F(x_\mu)=-\Pi_{\mathcal V}c/\mu$, are standard
self-concordance tools
\cite{nesterov1994-interior-point-polynomial-algorithms-convex,renegar2001-mathematical-view-ipm}.
The contribution here is their universal-over-barriers conditioning statement
and its geometric consequences, not new containment machinery.  In
particular, the evidence for novelty is a prior-art audit, not a claim that
these proof ingredients are new.
\end{remark}

\subsection{Sublevel chords and two containment facts}
\label{sec:conditioning-setup}

Let $P\subset\R^N$ be compact and convex, with affine hull
$x^0+\mathcal V$ and nonempty relative interior $P^\circ$.  Let $W$ have
orthonormal columns spanning the $d$-dimensional tangent space $\mathcal V$.
We minimize the nonconstant linear functional $c^\top x$, so
$c_{\mathcal V}:=\Pi_{\mathcal V}c\ne0$.  Put
\begin{equation}
 \mathrm{OPT}=\min_{y\in P}c^\top y,
 \qquad
 \Phi=\{y\in P:c^\top y=\mathrm{OPT}\}.
 \label{eq:conditioning-optimal-face}
\end{equation}
For $t\geq0$, define the objective sublevel set, its diameter, and the longest
sublevel chord from an interior point by
\begin{equation}
 L(t)=\{y\in P:c^\top y\leq\mathrm{OPT}+t\},\qquad
 D(t)=\operatorname{diam} L(t),\qquad
 \ell(x)=\sup_{y\in L(c^\top x-\mathrm{OPT})}\norm{y-x}_2.
 \label{eq:sublevel-chord-definitions}
\end{equation}
Thus $D$ is nondecreasing.  If $x$ has objective gap $g$, then $x\in L(g)$
and the triangle inequality gives
\begin{equation}
 \frac{D(g)}2\leq\ell(x)\leq D(g).
 \label{eq:diameter-chord-comparison}
\end{equation}
Indeed, for any $y,z\in L(g)$, at least one of
$\norm{x-y}_2,\norm{x-z}_2$ is at least $\norm{y-z}_2/2$.

Let $F$ be a $\nu_F$-self-concordant barrier for $P$ in the sense of
\cref{def:self-concordant-barrier}.  We use two standard consequences of
self-concordance.  For $x\in P^\circ$,
\begin{align}
 y\in x^0+\mathcal V,\quad
 \norm{y-x}_x<1&\quad\Longrightarrow\quad y\in P^\circ,
 \tag{Dikin}\label{eq:dikin-containment}\\
 y\in P,\quad
 \inner{\nabla F(x)}{y-x}\geq0
 &\quad\Longrightarrow\quad
 \norm{y-x}_x\leq C_F,
 \qquad C_F:=\nu_F+2\sqrt{\nu_F}.
 \tag{asymmetric}\label{eq:asymmetric-containment}
\end{align}
The first implication is understood for points in the affine hull; its
contrapositive says that a boundary or exterior point has local distance at
least one.  We also use semiboundedness,
\begin{equation}
 \inner{\nabla F(x)}{y-x}\leq\nu_F\qquad(y\in P).
 \label{eq:barrier-semiboundedness}
\end{equation}
These facts follow from the standard self-concordant-barrier containment
theorems; affine restriction preserves them and does not increase the
declared parameter
\cite{nesterov1994-interior-point-polynomial-algorithms-convex,renegar2001-mathematical-view-ipm}.

At an exact central point $x=x_F(\mu)$, first-order optimality on the affine
slice is
\begin{equation}
 W^\top(c+\mu\nabla F(x))=0.
 \label{eq:centrality-tangent}
\end{equation}
Consequently, for every $y\in P$,
\begin{equation}
 c^\top(x-y)=\mu\inner{\nabla F(x)}{y-x}\leq\nu_F\mu.
 \label{eq:central-objective-bound-pointwise}
\end{equation}
Taking $y\in\Phi$ proves the convention required in
\cref{rem:g-convention}:
\begin{equation}
 0<g_F(\mu)\leq\nu_F\mu.
 \label{eq:central-gap-bound}
\end{equation}
The strict inequality holds because a relative-interior minimizer of a
nonconstant linear functional cannot exist.  Notice that
\eqref{eq:central-gap-bound} is generally an inequality.  The equality
$\nu\mu$ in \cref{def:central-path} is the canonical \emph{primal--dual}
gap, not necessarily the primal objective gap $g$ used here.

\subsection{Rigidity of the Hessian spectrum}
\label{sec:conditioning-rigidity}

We begin with the barrier-independent lower edge.  It does not require
centrality.

\begin{theorem}[Universal minimum-eigenvalue rigidity]
\label{thm:lambda-min-rigidity}
Under the setup of \cref{sec:conditioning-setup}, for every
$\nu_F$-self-concordant barrier $F$ for $P$ and every $x\in P^\circ$,
\begin{equation}
 \lambda_{\min}\!\left(W^\top\nabla^2F(x)W\right)
 \geq \frac1{\ell(x)^2}.
 \label{eq:lambda-min-rigidity}
\end{equation}
Here $\ell(x)$ is formed from the objective level of that same point, whether
or not $x$ is central.
\end{theorem}

\begin{proof}
Fix a Euclidean unit vector $h\in\mathcal V$ and replace it by $-h$ if
necessary so that $c^\top h\leq0$.  Compactness implies that the ray
$x+th$ leaves $P^\circ$ in finite time.  More precisely, if
\[
 \tau=\sup\{t\geq0:x+sh\in P^\circ\text{ for }0\leq s<t\},
\]
then $0<\tau<\infty$ and $y=x+\tau h\in P$.  The objective does not
increase along the ray, so $y\in L(c^\top x-\mathrm{OPT})$ and hence
$\tau\leq\ell(x)$.  Since $y\notin P^\circ$, the contrapositive of
\eqref{eq:dikin-containment} gives
\[
 1\leq\norm{\tau h}_x^2
   =\tau^2\,\inner{\nabla^2F(x)h}{h}.
\]
Thus the Rayleigh quotient is at least
$\tau^{-2}\geq\ell(x)^{-2}$.  The Hessian form is even in $h$, so the
initial sign choice loses no direction.  Minimizing over the Euclidean unit
sphere in $\mathcal V$ proves \eqref{eq:lambda-min-rigidity}.
\end{proof}

The opposite spectral edge is forced by the distance in objective value to
the optimum.  Combining it with \cref{thm:lambda-min-rigidity} gives the
sublevel chord law.

\begin{theorem}[Sublevel chord law]
\label{thm:sublevel-chord-law}
Under the setup of \cref{sec:conditioning-setup}, let
$x=x_F(\mu)$ be an exact central point of any
$\nu_F$-self-concordant barrier, and set $g=g_F(\mu)$.  Then
\begin{align}
 \lambda_{\max}\!\left(W^\top\nabla^2F(x)W\right)
 &\geq \frac{\norm{c_{\mathcal V}}_2^2}{g^2},
 \label{eq:chord-lmax}\\
 \frac1{\ell(x)^2}
 \leq\lambda_{\min}\!\left(W^\top\nabla^2F(x)W\right)
 &\leq\frac{C_F^2}{\ell(x)^2},
 \label{eq:chord-lmin}\\
 \kappa_{\rm red}(\mu;F)
 &\geq
 \left(\frac{\ell(x)\norm{c_{\mathcal V}}_2}{C_Fg}\right)^2
 \geq
 \left(\frac{D(g)\norm{c_{\mathcal V}}_2}{2C_Fg}\right)^2.
 \label{eq:sublevel-chord-floor}
\end{align}
In particular, using only the $\mu$ parametrization,
\begin{equation}
 \kappa_{\rm red}(\mu;F)
 \geq
 \left(
  \frac{D(g_F(\mu))\norm{c_{\mathcal V}}_2}
       {2C_F\nu_F\mu}
 \right)^2.
 \label{eq:sublevel-chord-floor-mu}
\end{equation}
There is no small-$\mu$ threshold and no polyhedrality, nondegeneracy, or
strict-complementarity hypothesis.
\end{theorem}

\begin{proof}
The gap statement \eqref{eq:central-gap-bound} was proved above.  For the
largest eigenvalue, let
\[
 \widehat h=-\frac{c_{\mathcal V}}{\norm{c_{\mathcal V}}_2},
 \qquad
 t_* =\frac{g}{\norm{c_{\mathcal V}}_2}.
\]
The affine point $z=x+t_*\widehat h$ has $c^\top z=\mathrm{OPT}$.  If it
were in $P^\circ$, it would be a relative-interior minimizer of a linear
functional, forcing $c_{\mathcal V}=0$.  Therefore $z\notin P^\circ$, and
\eqref{eq:dikin-containment} gives
\[
 \inner{\nabla^2F(x)\widehat h}{\widehat h}
 \geq t_*^{-2}=\frac{\norm{c_{\mathcal V}}_2^2}{g^2}.
\]
This proves \eqref{eq:chord-lmax}.

For the upper bound in \eqref{eq:chord-lmin}, take any $y\in L(g)$.  By
\eqref{eq:centrality-tangent},
\[
 \inner{\nabla F(x)}{y-x}
 =-\frac{c^\top(y-x)}\mu\geq0.
\]
Asymmetric containment \eqref{eq:asymmetric-containment} therefore gives
$\norm{y-x}_x\leq C_F$.  If $y\ne x$, the Euclidean unit vector
$v=(y-x)/\norm{y-x}_2$ satisfies
\[
 \inner{\nabla^2F(x)v}{v}
 \leq\frac{C_F^2}{\norm{y-x}_2^2}.
\]
The smallest eigenvalue is no larger than this Rayleigh quotient.  Taking a
sequence of $y$ whose distance from $x$ tends to $\ell(x)$ proves the upper
bound; the lower bound is \cref{thm:lambda-min-rigidity}.  Dividing
\eqref{eq:chord-lmax} by the upper bound in \eqref{eq:chord-lmin}, then using
\eqref{eq:diameter-chord-comparison}, proves
\eqref{eq:sublevel-chord-floor}.  Finally $g\leq\nu_F\mu$ proves
\eqref{eq:sublevel-chord-floor-mu}.
\end{proof}

The maximum-eigenvalue constant in \eqref{eq:chord-lmax} can be exact.  On
the simplex problem $\min x_3$ subject to $x_1+x_2+x_3=1$, numerical central
paths for the two tested barriers give
$\lambda_{\max}g^2/\norm{c_{\mathcal V}}_2^2\to1$; for the ordinary log
barrier, $g/\mu\to1$ and
$\lambda_{\max}\mu^2\to2/3=\norm{c_{\mathcal V}}_2^2$.

The chord idea controls more than the two extreme eigenvalues.  For an
unoriented one-dimensional subspace $[h]\subset\mathcal V$, choose its unit
representative $\widehat h$ with $c^\top\widehat h<0$.  If
$c^\top h=0$, choose the sign having the longer exit time, breaking a tie
arbitrarily.  Define
\begin{equation}
 \tau_x([h])=\sup\{t\geq0:x+s\widehat h\in P^\circ
                     \text{ for }0\leq s<t\}.
 \label{eq:directional-exit-time}
\end{equation}
The selected exit point is a sublevel point, so
$0<\tau_x([h])\leq\ell(x)$.  The equatorial convention makes $\tau_x$
well-defined on unoriented lines; continuity is not needed.

For $j=1,\ldots,d$, define the two width profiles
\begin{align}
 w_j^-(x)
 &=\inf_{\substack{U\subseteq\mathcal V\\\operatorname{codim}U=j-1}}
   \ \sup_{\substack{h\in U\\\norm h_2=1}}\tau_x([h]),
 \label{eq:lower-width-profile}\\
 w_j^+(x)
 &=\sup_{\substack{S\subseteq\mathcal V\\\dim S=j}}
   \ \inf_{\substack{h\in S\\\norm h_2=1}}\tau_x([h]).
 \label{eq:upper-width-profile}
\end{align}
We use infima and suprema because the sign choice can make the exit-time
function discontinuous on $c^\perp$.

\begin{theorem}[Width-spectrum rigidity]
\label{thm:width-spectrum-rigidity}
In the setting of \cref{thm:sublevel-chord-law}, list the eigenvalues of
$W^\top\nabla^2F(x)W$ in nondecreasing order.  For every
$j=1,\ldots,d$,
\begin{equation}
 \frac1{(w_j^-(x))^2}
 \leq\lambda_j\!\left(W^\top\nabla^2F(x)W\right)
 \leq\frac{C_F^2}{(w_j^+(x))^2}.
 \label{eq:width-spectrum-sandwich}
\end{equation}
Moreover $w_j^+(x)\leq w_j^-(x)$ for every $j$, and
$w_1^+(x)=\sup_{[h]}\tau_x([h])\geq\ell(x)$.
\end{theorem}

\begin{proof}
For every unit $h$, Dikin containment at the exit point gives
\begin{equation}
 \inner{\nabla^2F(x)h}{h}\geq\tau_x([h])^{-2}.
 \label{eq:width-pointwise-lower}
\end{equation}
For the selected orientation, centrality gives
\[
 \inner{\nabla F(x)}{\tau_x([h])\widehat h}
 =-\frac{\tau_x([h])c^\top\widehat h}{\mu}\geq0.
\]
Asymmetric containment \eqref{eq:asymmetric-containment} applies directly
to the exit point and gives
\begin{equation}
 \inner{\nabla^2F(x)h}{h}\leq C_F^2\tau_x([h])^{-2}.
 \label{eq:width-pointwise-upper}
\end{equation}
The Rayleigh quotient is unchanged by orientation.

The two Courant--Fischer formulas for eigenvalues in nondecreasing order are
\begin{align*}
 \lambda_j&=\sup_{\operatorname{codim}U=j-1}
             \ \inf_{h\in U,\norm h_2=1}R(h),\\
 \lambda_j&=\inf_{\dim S=j}
             \ \sup_{h\in S,\norm h_2=1}R(h),
\end{align*}
where $R(h)=\inner{\nabla^2F(x)h}{h}$.  Applying
\eqref{eq:width-pointwise-lower} in the first formula yields
\[
 \lambda_j\geq
 \sup_U\frac1{(\sup_{h\in U}\tau_x([h]))^2}
 =\frac1{(w_j^-(x))^2}.
\]
Applying \eqref{eq:width-pointwise-upper} in the second formula gives
\[
 \lambda_j\leq C_F^2
 \inf_S\frac1{(\inf_{h\in S}\tau_x([h]))^2}
 =\frac{C_F^2}{(w_j^+(x))^2}.
\]

For any $j$-dimensional $S$ and codimension-$(j-1)$ subspace $U$, the
dimension formula gives a nonzero line $[h]\subseteq S\cap U$.  Hence
\[
 \inf_{v\in S,\norm v=1}\tau_x([v])
 \leq\tau_x([h])
 \leq\sup_{v\in U,\norm v=1}\tau_x([v]).
\]
Taking the supremum over $S$ and infimum over $U$ proves
$w_j^+\leq w_j^-$.  Finally, if $y\in L(g)$, the segment $[x,y)$ lies in
$P^\circ$.  The line through $y-x$ is oriented toward $y$ when the objective
strictly decreases, and our longer-exit convention gives the same conclusion
when it is unchanged.  Thus
$\tau_x([y-x])\geq\norm{y-x}_2$.  Taking the supremum over $y$ proves
$w_1^+\geq\ell(x)$.
\end{proof}

Thus even the interior spectrum is constrained by directional sublevel
widths, up to the explicit $C_F^2$ window and the gap between the two minimax
widths.  The empirical consequences for spectral clustering are deferred to
the next section.

\subsection{Canonical barriers attain the floor}
\label{sec:conditioning-achievability}

We next specialize to compact affine slices of the nonnegative orthant and of
the product PSD cone in \cref{sec:prelim-conic}.  To distinguish parameters
when two barriers are compared, write $\nu_{\rm can}=n$ for the LP log
barrier and $\nu_{\rm can}=\sum_jd_j$ for the product log-det barrier.  This
is the cone parameter denoted by $\nu$ in \cref{sec:prelim-conic}.

\begin{theorem}[Canonical achievability and gap-matched comparison]
\label{thm:canonical-achievability}
Consider either of the following strictly primal--dual feasible problems with
a compact primal feasible region and nonconstant objective.
\begin{enumerate}[label=\textup{(\roman*)}]
 \item An LP in \eqref{eq:lp-pair}, with canonical path
 $(x(\mu),y(\mu),s(\mu))$ and $M(\mu)=\norm{s(\mu)}_\infty$.
 \item A product-form SDP in \eqref{eq:sdp-pair}.  Its canonical path is
 $(X(\mu),y(\mu),S(\mu))$, and
 $M(\mu)=\max_j\norm{S_j(\mu)}_2$.
\end{enumerate}
Fix $\bar\mu>0$ and suppose the tail constant
\begin{equation}
 M_{\rm tail}=\sup_{0<\mu\leq\bar\mu}M(\mu)
 \label{eq:dual-tail-constant}
\end{equation}
is finite.  Use the Euclidean metric for the LP and the product Frobenius
metric for the SDP.  Write
$g_{\rm can}(\bar\mu)=g_{F_{\rm can}}(\bar\mu)$ for the canonical primal
objective gap at the endpoint of the tail.  The notation
$\kappa_{\rm red}^{\rm can}(g)$ and $\kappa_{\rm red}^{F}(g)$ means the
reduced condition number at the central point of gap $g$ for the named
barrier, using the $g$-parametrization in \cref{def:central-path} on an
interval where the named objective-gap map $g_F$ is one-to-one.  For a
canonical central point with $0<\mu\leq\bar\mu$ and primal objective gap $g$,
\begin{align}
 \lambda_{\min}(H_{\rm red})&\geq\ell(x)^{-2},
 \label{eq:canonical-lmin}\\
 \lambda_{\max}(H_{\rm red})&\leq M_{\rm tail}^2\mu^{-2}
 \leq M_{\rm tail}^2\nu_{\rm can}^2g^{-2},
 \label{eq:canonical-lmax}\\
 \kappa_{\rm red}^{\rm can}(g)&\leq
 \left(\frac{M_{\rm tail}\nu_{\rm can}D(g)}g\right)^2.
 \label{eq:canonical-kappa-upper}
\end{align}
For every gap $0<g\leq g_{\rm can}(\bar\mu)$ attained on the canonical tail,
any $\nu_F$-self-concordant barrier $F$ on the same feasible region, and any
exact $F$-central point having that gap,
\begin{equation}
 \kappa_{\rm red}^{\rm can}(g)
 \leq
 \left(
  \frac{2M_{\rm tail}\nu_{\rm can}
              (\nu_F+2\sqrt{\nu_F})}
       {\norm{c_{\mathcal V}}_2}
 \right)^2
 \kappa_{\rm red}^{F}(g).
 \label{eq:gap-matched-barrier-comparison}
\end{equation}
The comparison is between points of equal objective gap; their values of
$\mu$ need not agree.
\end{theorem}

\begin{proof}
The first bound is \cref{thm:lambda-min-rigidity}.  For the LP, canonical
complementarity gives $x_i=\mu/s_i$, and therefore, for every Euclidean unit
tangent vector $h$,
\[
 \inner{\nabla^2F_{\log}(x)h}{h}
 =\sum_i\frac{h_i^2}{x_i^2}
 \leq\frac{M_{\rm tail}^2}{\mu^2}.
\]
For the SDP, $X_j^{-1}=S_j/\mu$, and for a unit product-Frobenius tangent
direction $H$,
\begin{align*}
 D^2F_{\log\det}(X)[H,H]
 &=\sum_j\norm{X_j^{-1/2}H_jX_j^{-1/2}}_F^2\\
 &\leq \max_j\norm{X_j^{-1}}_2^2\sum_j\norm{H_j}_F^2
 \leq\frac{M_{\rm tail}^2}{\mu^2}.
\end{align*}
This proves the first inequality in \eqref{eq:canonical-lmax}.  The canonical
primal--dual gap is $\nu_{\rm can}\mu$, whereas its primal objective gap
satisfies $g\leq\nu_{\rm can}\mu$ by \eqref{eq:central-gap-bound}; hence
$\mu^{-1}\leq\nu_{\rm can}/g$.  Combining the largest-eigenvalue bound with
\eqref{eq:canonical-lmin} and $\ell(x)\leq D(g)$ proves
\eqref{eq:canonical-kappa-upper}.  Finally divide that upper bound by the
universal lower bound \eqref{eq:sublevel-chord-floor} for $F$ to obtain
\eqref{eq:gap-matched-barrier-comparison}.
\end{proof}

The tail boundedness in \eqref{eq:dual-tail-constant} is an explicit
instance hypothesis; it holds in the standard convergent canonical-path
setting.  The restriction to a tail is necessary because the corresponding
supremum can diverge as $\mu\to\infty$.  More importantly,
\eqref{eq:gap-matched-barrier-comparison} says that no alternative barrier can
improve the canonical reduced-Hessian condition number by more than a
$g$-independent instance factor.  It does not say that an arbitrary barrier
cannot be worse.

\subsection{A dictionary from error bounds to conditioning}
\label{sec:conditioning-dictionary}

The diameter $D(g)$ is the relevant sharpness modulus.  If, for some
$\theta\in[0,1]$ and $\gamma>0$,
\begin{equation}
 \operatorname{dist}(y,\Phi)\leq\gamma^{-1}
       (c^\top y-\mathrm{OPT})^\theta\qquad(y\in P),
 \label{eq:holder-error-bound}
\end{equation}
then choosing nearest optimal points for two members of $L(g)$ gives
\begin{equation}
 D(g)\leq\operatorname{diam}\Phi+2\gamma^{-1}g^\theta.
 \label{eq:diameter-from-error-bound}
\end{equation}
Conversely, $\Phi\subseteq L(g)$ and a central point $x$ of gap $g$ belongs
to $L(g)$.  For every $x^*\in\Phi$,
$g=c_{\mathcal V}^\top(x-x^*)\leq
\norm{c_{\mathcal V}}_2\norm{x-x^*}_2$, so
\begin{equation}
 D(g)\geq\operatorname{diam}\Phi,
 \qquad
 D(g)\geq\frac g{\norm{c_{\mathcal V}}_2}
 \quad\text{at every attained central gap }g.
 \label{eq:diameter-lower-basics}
\end{equation}

\begin{corollary}[Conditioning regimes]
\label{cor:conditioning-regimes}
Assume the hypotheses of \cref{thm:sublevel-chord-law}; for canonical upper
bounds also assume \cref{thm:canonical-achievability}, and restrict those
conclusions to the canonical tail.
\begin{enumerate}[label=\textup{(\roman*)}]
 \item If $\dim\Phi\geq1$, then every fixed $\nu_F$-self-concordant barrier
 satisfies $\kappa_{\rm red}^{F}(g)=\Omega(g^{-2})$.  The canonical barrier
 satisfies $\kappa_{\rm red}^{\rm can}(g)=\Theta(g^{-2})$.
 \item If $P$ is a polytope with a unique optimum, including a degenerate
 vertex, then $D(g)=\Theta(g)$ as $g\downarrow0$.  Consequently the canonical
 log barrier has $\kappa_{\rm red}^{\log}(g)=\Theta(1)$.
 \item If the optimum is unique and $D(g)=\Theta(\sqrt g)$, as under a
 two-sided quadratic-growth geometry, every fixed barrier has
 $\kappa_{\rm red}^{F}(g)=\Omega(g^{-1})$ and the canonical barrier has
 $\kappa_{\rm red}^{\rm can}(g)=\Theta(g^{-1})$.
\end{enumerate}
The constants may depend on the instance and, for a noncanonical barrier, on
$\nu_F$.  Whenever the canonical objective gap is $\Theta(\mu)$, as in the
standard strictly complementary tail, the same three canonical rates are
$\Theta(\mu^{-2})$, $\Theta(1)$, and $\Theta(\mu^{-1})$, respectively.
\end{corollary}

\begin{proof}
If $\dim\Phi\geq1$, the face contains two distinct points and hence
$\operatorname{diam}\Phi>0$; compactness makes this diameter finite and
attained.
Substitution of $D(g)\geq\operatorname{diam}\Phi$ in
\eqref{eq:sublevel-chord-floor} proves the universal lower bound.  Since
$D(g)\leq\operatorname{diam}P$, \eqref{eq:canonical-kappa-upper} gives the
matching canonical upper bound.

For a polytope, the solution set of a feasible bounded linear program is a
weak sharp minimum
\cite{mangasarian1979-nonlinear-perturbation-linear-programs,burke1993-weak-sharp-minima}.
With a unique optimum $x^*$, this gives \eqref{eq:holder-error-bound} with
$\theta=1$, hence $D(g)\leq2g/\gamma$.  For the reverse bound, choose a
feasible $z\ne x^*$ and set
$\Delta=c^\top z-\mathrm{OPT}>0$.  For every $0<g\leq\Delta$, convexity makes
\[
 y_g=\left(1-\frac g\Delta\right)x^*+\frac g\Delta z
\]
feasible, and $c^\top y_g=\mathrm{OPT}+g$.  Thus
\[
 D(g)\geq\norm{y_g-x^*}_2
 =\frac{\norm{z-x^*}_2}{\Delta}\,g.
\]
Combining these two bounds with
\cref{thm:sublevel-chord-law,thm:canonical-achievability} proves (ii).
Part (iii) follows by substituting
$D(g)=\Theta(\sqrt g)$ in the same two theorems.  The final conversion is
immediate under the separately stated relation $g=\Theta(\mu)$.
\end{proof}

\begin{remark}
The statement uniform over barriers in \cref{cor:conditioning-regimes} is the
lower rate.  The matching upper rate is established for the canonical
barrier.  Thus the barrier-optimal rate in the positive-dimensional-face
case is $\Theta(g^{-2})$; the results do not supply an upper bound of that
order for every conceivable self-concordant barrier.
\end{remark}

Part (ii) includes degenerate vertices: nondegeneracy is not needed for weak
sharpness.  Part (iii) explains why the reduced primal log-det Hessian at a
generic curved, unique SDP optimum can have condition number
$\Theta(\mu^{-1})$, even though its largest eigenvalue and a primal--dual
Schur complement can live at the $\mu^{-2}$ scale.  This is the conic analogue
of the $\mu^{-1}$ structured logarithmic-barrier behavior identified for
nonlinear programs by
\citet{wright1994-properties-hessian-logarithmic-barrier}.

The same substitution produces fractional exponents.  Sturm proves an
$O(\epsilon^{2^{-d}})$ distance bound when the relevant semidefinite
feasibility system has singularity degree $d$
\cite{sturm2000-error-bounds-linear-matrix-inequalities}.  An upper error
bound alone does not force a conditioning lower bound: its exponent must also
be attained by sublevel chords.  When it is sharp, the result is immediate.

\begin{corollary}[Attained H\"older and singularity-degree laws]
\label{cor:fractional-conditioning-laws}
Suppose that, along central gaps $g\downarrow0$,
$D(g)\geq a g^\theta$ for constants $a>0$ and $\theta\in[0,1]$.  Then every
$\nu_F$-self-concordant barrier satisfies
\begin{equation}
 \kappa_{\rm red}^{F}(g)
 \geq
 \frac{a^2\norm{c_{\mathcal V}}_2^2}{4C_F^2}
 g^{2\theta-2}.
 \label{eq:holder-conditioning-law}
\end{equation}
If, in addition, the problem is a canonical LP or product-form SDP satisfying
all hypotheses of \cref{thm:canonical-achievability} and
$D(g)=O(g^\theta)$, then the canonical barrier attains
$\Theta(g^{2\theta-2})$ on its tail.  In particular, when the Sturm exponent
$\theta=2^{-d}$ is attained, the exponents are
\begin{equation}
 d=0:\ 0,\qquad d=1:\ -1,\qquad d=2:\ -\tfrac32,
 \qquad
 2\theta-2=2^{1-d}-2\longrightarrow-2.
 \label{eq:singularity-degree-exponents}
\end{equation}
\end{corollary}

\begin{proof}
Insert $D(g)\geq ag^\theta$ into
\eqref{eq:sublevel-chord-floor}.  If the matching upper bound on $D$ holds,
combine the resulting lower bound with
\eqref{eq:canonical-kappa-upper}.  Substituting $\theta=2^{-d}$ gives
\eqref{eq:singularity-degree-exponents}.
\end{proof}

The following $3\times3$ problem shows that the $d=2$ fractional law occurs
in a concrete SDP rather than only in the error-bound formalism.

\begin{proposition}[An explicit $g^{-3/2}$ witness]
\label{prop:degree-two-witness}
Consider
\begin{equation}
 \begin{aligned}
  \min_{X\in\Sym^3}\quad &X_{22}\\
  \text{subject to}\quad &\tr X=1,\qquad X_{33}=X_{12},
  \qquad X\succeq0.
 \end{aligned}
 \label{eq:degree-two-sdp}
\end{equation}
Its unique optimum is $X^*=e_1e_1^\top$, its optimality system reaches the
minimal face in two facial-reduction steps, and
\begin{equation}
 D(g)=\Theta(g^{1/4}).
 \label{eq:degree-two-diameter}
\end{equation}
Consequently every self-concordant barrier obeys
$\kappa_{\rm red}(g)=\Omega(g^{-3/2})$, while the log-det barrier, under the
hypotheses of \cref{thm:canonical-achievability}, has
$\kappa_{\rm red}(g)=\Theta(g^{-3/2})$ on its canonical tail.
\end{proposition}

\begin{proof}
Write a feasible matrix as
\[
 X=\begin{bmatrix}
      a&b&t\\ b&d&e\\ t&e&b
    \end{bmatrix},
 \qquad a+d+b=1,\qquad X\succeq0.
\]
Here $d=X_{22}$ is the objective and $b=X_{12}=X_{33}\geq0$.  If $d=0$,
positive semidefiniteness makes the second row and column zero, hence $b=e=0$.
Then the third diagonal entry is zero, so the third row and column vanish and
$t=0$; the trace constraint gives $a=1$.  The optimum is therefore unique.

For $X\in L(g)$, the $2\times2$ principal minors give
\[
 b^2\leq ad\leq d\leq g,
 \qquad
 t^2\leq ab\leq b\leq\sqrt g,
 \qquad
 e^2\leq db\leq g^{3/2}.
\]
Also $|a-1|=d+b\leq g+\sqrt g$.  Thus
$\norm{X-X^*}_F=O(g^{1/4})$, and the triangle inequality gives
$D(g)=O(g^{1/4})$.

For the reverse bound, let $0<g\leq1/16$, set
\[
 d=g,\qquad b=\frac{\sqrt g}{2},\qquad
 a=1-g-\frac{\sqrt g}{2},\qquad e=0,\qquad
 t=\pm\frac{g^{1/4}}{2\sqrt2}.
\]
Then $a\geq13/16$.  The one- and two-dimensional principal minors are
positive, and
\[
 \det X
 =ag b-b^3-gt^2
 =g^{3/2}\left(\frac a2-\frac14\right)>0.
\]
Both sign choices therefore give positive-definite feasible points in
$L(g)$.  Their Frobenius distance is
$2\sqrt2|t|=g^{1/4}$, proving the lower bound in
\eqref{eq:degree-two-diameter}.

For completeness, the two facial reductions can be read directly from the
optimality equations.  The first exposing equation $X_{22}=0$ removes the
$e_2$ row and column; then $X_{12}=0$ turns the affine equation
$X_{33}=X_{12}$ into $X_{33}=0$, which removes the $e_3$ row and column.
No one-step positive-semidefinite exposing combination can do both: a PSD
linear combination of the trace, $X_{33}-X_{12}$, and $X_{22}$ equations
whose $(1,1)$ entry is zero must have zero first row, forcing the coefficient
of $X_{33}-X_{12}$ to vanish.  Thus two steps are necessary and sufficient.
Finally apply \cref{cor:fractional-conditioning-laws} with $\theta=1/4$.
\end{proof}

For the log-det barrier, the leading constant can also be computed
exactly. For sufficiently small positive gap $g$, the center and central parameter are
\[
 X_{\log}(g)=\begin{pmatrix}a&b&0\\b&g&0\\0&0&b\end{pmatrix},\qquad
 b=\frac{-g+\sqrt{3g-2g^2}}3,\quad a=1-g-b,\quad
 \mu=\frac{ag-b^2}{1-b-2g}.
\]
Indeed, sign symmetry forces the two displayed zero entries, and
stationarity of $g/\mu-\log b-\log(g(1-g-b)-b^2)$ gives
$3b^2+2gb+g^2-g=0$ and the stated $\mu$. Strict convexity makes
this interior stationary point the unique center. The spectral calculation
in \cref{sec:spectra} gives
\[
 \kappa_{\rm red}\mu^{3/2}\longrightarrow\frac{2\sqrt2}{7},
 \qquad
 \kappa_{\rm red}g^{3/2}\longrightarrow\frac{3\sqrt3}{7}.
\]
The first constant is approximately $0.40406$, refining the earlier
numerical value $0.40$. The exact center, Frobenius-metric spectrum, and
these limits are verified in Lean; the accompanying report
\texttt{formal/FRACTIONAL\_SDP.md} states the verification scope.
The geometric and singularity-degree proof above is separate from that
formalization.

Near-optimal faces yield a related finite plateau rather than a divergent
tail.

\begin{corollary}[Near-degeneracy plateau]
\label{cor:near-degeneracy-plateau}
Suppose a family of compact instances has a scale $\vartheta>0$ and constants
independent of $g,\vartheta$ such that
\begin{equation}
 D(g)=
 \begin{cases}
  \Theta(g/\vartheta),&0<g\lesssim\vartheta,\\
  \Theta(1),&\vartheta\lesssim g\leq g_0.
 \end{cases}
 \label{eq:near-face-diameter-law}
\end{equation}
Assume also that $\norm{c_{\mathcal V}}_2$ is bounded below and that the
barrier parameters under comparison are bounded above uniformly in
$\vartheta$; for the canonical upper bound, assume the same of
$M_{\rm tail}\nu_{\rm can}$ and assume
$g_0\leq g_{\rm can}(\bar\mu)$ uniformly in $\vartheta$.
Then every self-concordant barrier has the corresponding floor
\begin{equation}
 \kappa_{\rm red}^{F}(g)=
 \begin{cases}
  \Omega(\vartheta^{-2}),&0<g\lesssim\vartheta,\\
  \Omega(g^{-2}),&\vartheta\lesssim g\leq g_0,
 \end{cases}
 \label{eq:near-face-floor}
\end{equation}
up to its explicit $C_F$ and instance constants.  The canonical barrier
matches both orders on the tail under
\cref{thm:canonical-achievability}.  Thus the
$g^{-2}$ growth stops near $g\asymp\vartheta$ at height
$\Theta(\vartheta^{-2})$.
\end{corollary}

\begin{proof}
In the first range, $D(g)/g=\Theta(\vartheta^{-1})$; in the second it is
$\Theta(g^{-1})$.  Apply the lower bound
\eqref{eq:sublevel-chord-floor} and the canonical upper bound
\eqref{eq:canonical-kappa-upper}.
\end{proof}

For the simplex family
$\min\{x_3+\vartheta x_2:x_1+x_2+x_3=1,\ x\geq0\}$, the log-barrier
plateau has the following precise two-scale limit.  For each fixed
$\vartheta>0$,
\begin{equation}
 \lim_{\mu\downarrow0}\kappa_{\rm red}(\mu)
 =
 \frac{1+\vartheta^2+\sqrt{1-\vartheta^2+\vartheta^4}}
      {1+\vartheta^2-\sqrt{1-\vartheta^2+\vartheta^4}},
 \qquad
 \lim_{\vartheta\downarrow0}\vartheta^2
       \left(\lim_{\mu\downarrow0}\kappa_{\rm red}(\mu)\right)=\frac43.
 \label{eq:simplex-plateau-double-limit}
\end{equation}
To verify the inner limit, write the equality multiplier as
$\alpha_\mu=\mu/x_1$.  Centrality gives
$x_2=\mu/(\vartheta+\alpha_\mu)$,
$x_3=\mu/(1+\alpha_\mu)$, and $\alpha_\mu\to0$.  On the tangent plane, put
$h=(-u-v,u,v)$.  After multiplication by $\mu^2$, the limiting Rayleigh
quotient is
\[
 \frac{\vartheta^2u^2+v^2}{2(u^2+uv+v^2)}.
\]
Its generalized eigenvalues are
$\bigl(1+\vartheta^2\pm\sqrt{1-\vartheta^2+\vartheta^4}\bigr)/3$, whose
ratio is the first expression in
\eqref{eq:simplex-plateau-double-limit}; expansion at
$\vartheta=0$ gives the second limit.  The computations at
$\vartheta=10^{-2}$ and $10^{-3}$ round
$\vartheta^2\lim_{\mu\downarrow0}\kappa_{\rm red}(\mu)$ to $4/3$.

\subsection{Transfer to approximate centrality}
\label{sec:conditioning-approximate}

The floor is stable in the short-step neighborhood.  All norms in the next
statement are restricted to the tangent space.

\begin{corollary}[Approximate-centrality transfer]
\label{cor:approximate-centrality-transfer}
Fix $\mu>0$, let $x_*=x_F(\mu)$, and suppose $x\in P^\circ$ satisfies
\begin{equation}
 \lambda_\mu(x):=\norm{c/\mu+\nabla F(x)}_x^*\leq\rho\leq\frac16.
 \label{eq:approximate-centrality-decrement}
\end{equation}
Then, with $H_*=W^\top\nabla^2F(x_*)W$ and
$H=W^\top\nabla^2F(x)W$,
\begin{equation}
 \norm{x-x_*}_{x_*}\leq\frac{\rho}{(1-\rho)^2}\leq\frac14,
 \qquad
 \kappa(H)\geq\left(\frac34\right)^4\kappa(H_*).
 \label{eq:approximate-centrality-transfer}
\end{equation}
Consequently $x$ inherits the right-hand sides of
\eqref{eq:sublevel-chord-floor} and
\eqref{eq:sublevel-chord-floor-mu}, evaluated at the exact center $x_*$,
with the additional factor $(3/4)^4$.
\end{corollary}

\begin{proof}
Restrict $f=F+c^\top(\cdot)/\mu$ to the affine slice.  It is standard
self-concordant, has Hessian $\nabla^2F$, and is minimized at $x_*$.  With
\[
 \omega(t)=t-\log(1+t),\qquad
 \omega_*(t)=-t-\log(1-t),
\]
the self-concordant Newton estimates give
\[
 \omega(\norm{x-x_*}_{x_*})
 \leq f(x)-f(x_*)
 \leq\omega_*(\lambda_\mu(x))
 \leq\omega_*(\rho).
\]
The elementary scalar inequality
$\omega^{-1}(\omega_*(\rho))\leq\rho/(1-\rho)^2$ follows by substituting
$t=\rho/(1-\rho)^2$ and differentiating the difference
$\omega(t)-\omega_*(\rho)$ from its value zero at $\rho=0$.  This proves the
distance bound, and $\rho\leq1/6$ gives
$\rho/(1-\rho)^2\leq6/25<1/4$.

If $r=\norm{x-x_*}_{x_*}<1$, self-concordant Hessian comparison gives
\[
 (1-r)^2H_*\preceq H\preceq(1-r)^{-2}H_*.
\]
It follows that
$\lambda_{\max}(H)\geq(1-r)^2\lambda_{\max}(H_*)$ and
$\lambda_{\min}(H)\leq(1-r)^{-2}\lambda_{\min}(H_*)$; hence
$\kappa(H)\geq(1-r)^4\kappa(H_*)\geq(3/4)^4\kappa(H_*)$.
\end{proof}

\subsection{Implications and open geometric questions}
\label{sec:conditioning-implications}

For QIPMs that solve barrier-Newton systems, changing the self-concordant
barrier cannot remove a sublevel-chord conditioning obstruction.  At target
primal objective gap $\epsilon$, every barrier pays at least
\[
 \left(
  \frac{D(\epsilon)\norm{c_{\mathcal V}}_2}
       {2(\nu_F+2\sqrt{\nu_F})\epsilon}
 \right)^2
\]
in the reduced Hessian, and, for canonical LPs or product-form SDPs with
finite $M_{\rm tail}$ and $0<\epsilon\leq g_{\rm can}(\bar\mu)$ attained on
the canonical tail, the canonical barrier is within the explicit instance
factor in \eqref{eq:gap-matched-barrier-comparison}.  In particular,
positive-dimensional optimal faces impose a quadratic inverse-gap floor,
curved unique optima can impose a linear inverse-gap floor, and sharp
singularity-degree exponents can impose fractional laws.

This conclusion is deliberately scoped to $\kappa_{\rm red}$ as defined in
\cref{def:reduced-kappa}.  It makes no claim about condition-number-free QIPM
approaches such as
\citet{apers2026-quantum-speedups-linear-programming-interior}, about
preconditioners outside barrier-Hessian form, or about scalar congruences.
Those access-model and system-representation questions are treated in later
sections.

\begin{openproblem}[Degenerate-vertex constants]
\label{open:degenerate-vertex-constants}
For a polytope with a unique degenerate optimal vertex, characterize the
finite limiting constant (or limiting spectrum) of the canonical reduced log
Hessian directly from the active-set geometry.
\end{openproblem}

\begin{openproblem}[Facial characterization of chord growth]
\label{open:facial-chord-growth}
Characterize the exact growth of $D(g)$ for semidefinite programs from facial
data beyond the generic strict-complementarity and attained
singularity-degree cases.  In particular, determine when Sturm's
$2^{-d}$ exponent is sharp for objective sublevel diameters rather than only
an upper error-bound exponent.
\end{openproblem}
